\documentclass[a4paper,12pt]{article}
\usepackage[utf8]{inputenc}
\usepackage[german]{babel}
\usepackage{amsmath}
\usepackage{amssymb}
\usepackage{amsthm}
\usepackage{geometry}
\usepackage{graphicx}
\usepackage[hidelinks]{hyperref}
\usepackage{listings}
\usepackage{xcolor}
\usepackage{enumitem}
\usepackage{microtype}

\geometry{margin=2.5cm}

% Verhindert Überlauf über den rechten Rand
\setlength{\emergencystretch}{3em}
\sloppy

\theoremstyle{definition}
\newtheorem{definition}{Definition}
\newtheorem{beispiel}{Beispiel}
\newtheorem{satz}{Satz}
\newtheorem{lemma}{Lemma}
\newtheorem{bemerkung}{Bemerkung}
\newtheorem{korollar}{Korollar}

\setlist[itemize,1]{label=--}

% ---------- Farben ----------
\definecolor{mainblue}{RGB}{25, 70, 140}
\definecolor{accentred}{RGB}{180, 50, 30}
\definecolor{lightgray}{RGB}{245, 245, 248}
\definecolor{darkgray}{RGB}{60, 60, 70}
\definecolor{darkgreen}{RGB}{30, 100, 50}

\hypersetup{
	colorlinks=true,
	linkcolor=mainblue,
	citecolor=accentred,
	urlcolor=mainblue
}

%  Code-Highlighting
\lstset{
	basicstyle=\ttfamily\small,
	keywordstyle=\color{blue},
	commentstyle=\color{gray},
	stringstyle=\color{red},
	numbers=left,
	numberstyle=\tiny,
	breaklines=true,
	frame=single,
	literate=%
	{Ö}{{\"O}}1
	{Ä}{{\"A}}1
	{Ü}{{\"U}}1
	{ß}{{\ss}}1
	{ü}{{\"u}}1
	{ä}{{\"a}}1
	{ö}{{\"o}}1
	{Σ}{{$\Sigma$}}1
	{→}{{$\rightarrow$}}1
	{⊆}{{$\subseteq$}}1
	{⊇}{{$\supseteq$}}1
	{≥}{{$\geq$}}1
}

\title{Der geniale Karton: Nutzen, Beschaffenheit\\ und eine formale Analyse seiner strukturellen Eigenschaften}
\author{Stephan Epp}
\date{15. Juni 2026}

\begin{document}

\maketitle

\tableofcontents
\newpage

\section{Einführung}

Der Karton -- in seiner industriellen Ausprägung als Wellpappe-Verpackung -- ist eines der
unauffälligsten, zugleich aber wirkungsmächtigsten technischen Artefakte unserer
Zivilisation. Milliarden von Exemplaren durchlaufen täglich globale Logistikketten,
schützen empfindliche Güter vor Stoß, Druck und Feuchtigkeit, und werden nach Gebrauch
nahezu vollständig recycelt. Diese Arbeit untersucht den Karton nicht als triviales
Alltagsobjekt, sondern als Lösung zweier ineinandergreifender Optimierungsprobleme:
eines \emph{geometrischen} Problems (welche äußere Form minimiert den Materialeinsatz
bei vorgegebenem Volumen?) und eines \emph{strukturmechanischen} Problems (welche
innere Struktur des Wandmaterials maximiert die Tragfähigkeit bei minimalem
Materialeinsatz?).

\subsection{Problemstellung}

Gegeben sei ein zu verpackendes Gut mit Mindestvolumen $V > 0$. Gesucht ist ein
quaderförmiger Karton mit Kantenlängen $\ell, b, h > 0$ und $\ell b h \geq V$, dessen
Wandmaterial aus Wellpappe mit Decklagendicke $t_\ell$, Flötenhöhe $d$ und Flötenteilung
$\lambda$ besteht, sodass

\begin{itemize}
	\item der gesamte Materialeinsatz (Oberfläche $\times$ effektive Materialdicke,
	      gewichtet mit dem Take-up-Faktor der Wellenstruktur) minimal ist, und
	\item die resultierende Stapeldruckfestigkeit (\emph{Box Compression Test},
	      BCT) eine vorgegebene Mindestlast $F_{\min}$ nicht unterschreitet.
\end{itemize}

\textbf{Ergebnis dieser Arbeit:} Wir zeigen, dass die in der Praxis verwendeten
Kartonformen und Wellenprofile bereits sehr nahe an den formal hergeleiteten Optima
liegen -- der Karton ist, im wörtlichen Sinn, eine \emph{geniale} Konstruktion, weil er
zwei orthogonale Optimierungsprobleme (Makrogeometrie und Mikrostruktur) gleichzeitig
mit elementaren mathematischen Mitteln nahezu optimal löst.

\subsection{Motivation}

Die ökonomische und ökologische Bedeutung des Themas ist erheblich: Die
Verpackungsindustrie verarbeitet weltweit hunderte Millionen Tonnen Wellpappe pro Jahr.
Schon geringe relative Einsparungen im Materialeinsatz pro Karton -- etwa durch eine
Optimierung des Seitenverhältnisses oder der Wellenform um wenige Prozent -- summieren
sich auf gewaltige absolute Mengen an Papier, Energie und CO\textsubscript{2}.
Gleichzeitig ist der Karton ein dankbares Studienobjekt für die angewandte Mathematik:
Er verbindet klassische Extremwertaufgaben (Optimierung von Quadern unter
Volumenrestriktion), Differentialgeometrie (Bogenlängenintegrale sinusförmiger Profile),
Balkentheorie (Trägheitsmomente von Sandwich\-querschnitten, Euler-Knicken) und
empirische Ingenieurformeln (McKee-Formel) zu einem geschlossenen, formal
nachvollziehbaren Gesamtbild. Diese interdisziplinäre Synthese -- von der reinen
Extremwertrechnung bis zur empirisch validierten Ingenieurformel -- macht den Karton zu
einem idealen Beispiel dafür, wie rigorose formale Methoden auf scheinbar triviale
Alltagsobjekte angewendet werden können, um deren tatsächliche \emph{Genialität}
nachzuweisen.

\section{Grundlagen}

\subsection{Definitionen}

\begin{definition}[Karton]
	Ein Karton ist ein Quader $K = [0,\ell] \times [0,b] \times [0,h]$ mit
	Kantenlängen $\ell, b, h > 0$, dessen Begrenzungsflächen aus Wellpappe bestehen.
	Das Innenvolumen beträgt $V = \ell b h$, die gesamte Außenfläche (für einen
	geschlossenen Karton mit Deckel) beträgt
	\[
		S = 2(\ell b + b h + \ell h).
	\]
	Für einen offenen Karton (ohne Deckel) beträgt die Fläche
	$S = \ell b + 2h(\ell + b)$.
\end{definition}

\begin{definition}[Wellpappe]
	Wellpappe ist ein Sandwichmaterial, bestehend aus zwei ebenen Decklagen
	(\emph{liner}) der Dicke $t_\ell$ und einer dazwischenliegenden gewellten
	Mittellage (\emph{Flöte}) der Dicke $t_f$. Der Abstand zwischen den
	Mittelachsen der beiden Decklagen heißt Flötenhöhe $d$ (auch
	\emph{Kalliber} genannt).
\end{definition}

\begin{definition}[Flötenprofil]
	\label{def:flute}
	Das Profil der Mittellage einer Wellpappe wird im Sinusmodell beschrieben
	durch
	\[
		y(x) = A \sin\!\left(\frac{2\pi x}{\lambda}\right), \qquad x \in [0,\lambda],
	\]
	wobei $A = d/2$ die Amplitude (halbe Flötenhöhe) und $\lambda$ die
	\emph{Flötenteilung} (Periodenlänge, engl. \emph{pitch}) bezeichnet.
\end{definition}

\begin{definition}[Kantenstauchwiderstand (ECT) und Stapeldruckfestigkeit (BCT)]
	Der \emph{Edge Crush Test} (ECT) misst die Druckkraft pro Längeneinheit, die
	ein hochkant stehender Streifen Wellpappe parallel zu den Flöten aufnehmen
	kann, üblicherweise angegeben in $\mathrm{kN/m}$. Der \emph{Box Compression
	Test} (BCT) bezeichnet die maximale Druckkraft, die ein vollständiger,
	gefüllter Karton beim Stapeln (Belastung von oben) aufnehmen kann, angegeben
	in Newton.
\end{definition}

\begin{definition}[Knicklast]
	Die Knicklast (kritische Last) $P_{\mathrm{cr}}$ eines schlanken,
	plattenförmigen Bauteils ist die Druckkraft, ab der das Bauteil seine
	geradlinige Gleichgewichtslage verliert und seitlich ausweicht
	(\emph{Beulen}/\emph{Knicken}).
\end{definition}

\subsection{Grundidee}

Die Genialität des Kartons beruht auf der Entkopplung zweier Optimierungsebenen, die
unabhängig voneinander gelöst und anschließend kombiniert werden können:

\subsubsection{Makroebene: Form des Kartons}

Auf der Makroebene wird die äußere Quaderform $(\ell, b, h)$ so gewählt, dass bei
gegebenem Volumen $V$ die Oberfläche $S$ -- und damit der Materialverbrauch -- minimal
wird. Dies ist eine klassische Extremwertaufgabe mit Nebenbedingung.

\subsubsection{Mikroebene: Struktur des Wandmaterials}

Auf der Mikroebene wird das Wandmaterial selbst -- die Wellpappe -- so strukturiert,
dass bei gegebener Materialmasse (Papiergewicht pro Flächeneinheit) das
Flächenträgheitsmoment und damit die Biege- bzw. Knicksteifigkeit der Wand maximal
wird. Dies ist im Kern ein Problem der Balkentheorie.

\begin{beispiel}[Kombinatorische Explosion vs. analytische Lösung]
	Ein naiver Ansatz zur Optimierung eines Kartons würde diskrete Kombinationen
	aus Kantenlängen (in Millimeterschritten), Flötentypen (A, B, C, E, \ldots)
	und Decklagendicken durchprobieren. Bereits bei einer Diskretisierung von
	$\ell, b, h$ in $N = 1000$ Schritten und $4$ Flötentypen ergäben sich
	$N^3 \cdot 4 = 4 \cdot 10^9$ zu evaluierende Konfigurationen. Die in den
	folgenden Abschnitten herzuleitenden geschlossenen Formeln liefern die
	optimalen Parameter dagegen in $O(1)$ Rechenschritten -- ein Effizienzgewinn,
	der die Bedeutung formaler Analyse für ein scheinbar triviales Objekt
	unterstreicht.
\end{beispiel}

\section{Formale Modellierung}

\subsection{Geometrisches Modell}

Wir betrachten zunächst die Makroebene. Der Karton sei ein Quader mit Kantenlängen
$\ell, b, h > 0$ und Volumen $V = \ell b h$. Wir unterscheiden zwei Konstruktionsarten:

\begin{itemize}
	\item \textbf{Geschlossener Karton} (mit Deckel): Oberfläche
	      $S_{\mathrm{geschl}} = 2(\ell b + bh + \ell h)$.
	\item \textbf{Offener Karton} (ohne Deckel, z.\,B. Versandkarton mit
	      separatem Verschluss): Oberfläche
	      $S_{\mathrm{offen}} = \ell b + 2h(\ell + b)$.
\end{itemize}

\subsection{Optimierung der äußeren Form}

\begin{satz}[Optimalität des geschlossenen Kartons]
	\label{satz:wuerfel}
	Unter allen Quadern mit Kantenlängen $\ell, b, h > 0$ und festem Volumen
	$\ell b h = V$ minimiert der Würfel $\ell = b = h = V^{1/3}$ die Oberfläche
	$S_{\mathrm{geschl}} = 2(\ell b + bh + \ell h)$. Es gilt
	\[
		S_{\mathrm{geschl}} \geq 6V^{2/3},
	\]
	mit Gleichheit genau dann, wenn $\ell = b = h$.
\end{satz}

\begin{proof}
	Nach der Ungleichung vom arithmetischen und geometrischen Mittel (AM-GM) für
	die drei positiven Zahlen $\ell b$, $b h$, $\ell h$ gilt
	\[
		\frac{\ell b + bh + \ell h}{3} \;\geq\; \sqrt[3]{(\ell b)(bh)(\ell h)}
		= \sqrt[3]{\ell^2 b^2 h^2} = (\ell b h)^{2/3} = V^{2/3}.
	\]
	Multiplikation mit $3$ und anschließend mit $2$ liefert
	$S_{\mathrm{geschl}} = 2(\ell b + bh + \ell h) \geq 6V^{2/3}$. Gleichheit in
	der AM-GM-Ungleichung gilt genau dann, wenn alle drei Terme gleich sind, also
	$\ell b = bh = \ell h$. Aus $\ell b = bh$ folgt $\ell = h$ (für $b \neq 0$),
	und aus $bh = \ell h$ folgt $b = \ell$. Somit $\ell = b = h$, und wegen
	$\ell b h = V$ folgt $\ell = b = h = V^{1/3}$.
\end{proof}

\begin{satz}[Optimalität des offenen Kartons]
	\label{satz:offen}
	Unter allen offenen Kartons mit quadratischer Grundfläche $\ell = b$ und
	festem Volumen $V = \ell^2 h$ wird die Oberfläche
	$S_{\mathrm{offen}}(\ell) = \ell^2 + \dfrac{4V}{\ell}$ minimiert für
	\[
		\ell^\ast = (2V)^{1/3}, \qquad h^\ast = \frac{\ell^\ast}{2},
	\]
	mit Minimalwert $S_{\mathrm{offen}}^\ast = 3 \cdot 2^{2/3} V^{2/3} \approx
	4{,}7622 \, V^{2/3}$.
\end{satz}

\begin{proof}
	Mit $h = V/\ell^2$ ergibt sich
	$S_{\mathrm{offen}}(\ell) = \ell^2 + 2\ell \cdot \dfrac{2V}{\ell^2} = \ell^2 +
	\dfrac{4V}{\ell}$. Ableitung nach $\ell$:
	\[
		S_{\mathrm{offen}}'(\ell) = 2\ell - \frac{4V}{\ell^2}.
	\]
	Setzen wir $S_{\mathrm{offen}}'(\ell) = 0$, so folgt $2\ell^3 = 4V$, also
	$\ell^3 = 2V$ und $\ell^\ast = (2V)^{1/3}$. Die zweite Ableitung
	$S_{\mathrm{offen}}''(\ell) = 2 + \dfrac{8V}{\ell^3} > 0$ für alle
	$\ell > 0$ zeigt, dass es sich um ein (globales) Minimum handelt. Aus
	$h^\ast = V/(\ell^\ast)^2 = V/(2V)^{2/3} = (V/4)^{1/3}$ und
	$\ell^\ast = (2V)^{1/3} = 2 \cdot (V/4)^{1/3}$ (da
	$(2V)^{1/3} = 2^{1/3}V^{1/3}$ und $2 \cdot (V/4)^{1/3} = 2 \cdot
	V^{1/3}/4^{1/3} = 2^{1-2/3}V^{1/3} = 2^{1/3}V^{1/3}$, also identisch) folgt
	$\ell^\ast = 2h^\ast$, also $h^\ast = \ell^\ast/2$. Einsetzen liefert
	\[
		S_{\mathrm{offen}}^\ast = (\ell^\ast)^2 + \frac{4V}{\ell^\ast}
		= (2V)^{2/3} + \frac{4V}{(2V)^{1/3}}
		= 2^{2/3}V^{2/3} + 4V \cdot 2^{-1/3}V^{-1/3}
		= 2^{2/3}V^{2/3} + 2^{5/3}V^{2/3}.
	\]
	Wegen $2^{2/3} + 2^{5/3} = 2^{2/3}(1 + 2) = 3 \cdot 2^{2/3}$ folgt
	$S_{\mathrm{offen}}^\ast = 3 \cdot 2^{2/3}V^{2/3}$.
\end{proof}

\begin{bemerkung}[Skalierungsgesetz]
	Sowohl für den geschlossenen als auch für den offenen Karton skaliert die
	minimale Oberfläche mit $V^{2/3}$. Dies ist eine direkte Konsequenz der
	Dimensionsanalyse: Oberflächen besitzen die Dimension $[\text{Länge}]^2$,
	Volumina die Dimension $[\text{Länge}]^3$, und jede dimensionskonsistente
	Funktion $S(V)$ einer einzigen Variable $V$ muss daher von der Form
	$S(V) = c \cdot V^{2/3}$ sein. Die Konstante $c$ unterscheidet die beiden
	Konstruktionsarten ($c_{\mathrm{geschl}} = 6$, $c_{\mathrm{offen}} = 3 \cdot
	2^{2/3} \approx 4{,}7622$) und ist ein Maß für die \emph{Materialeffizienz}
	der jeweiligen Bauform. Bemerkenswert ist, dass der offene Karton -- obwohl
	er bei gleichem Volumen \emph{eine} Fläche (den Deckel) einspart --
	insgesamt nicht $1/6$, sondern nur etwa $20{,}6\,\%$ weniger Material
	benötigt, da sich die optimale Form gleichzeitig ändert (von einem Würfel zu
	einem Quader mit $\ell = b = 2h$).
\end{bemerkung}

\subsection{Mikrostruktur: Flötengeometrie}

\begin{lemma}[Take-up-Faktor der Wellpappe]
	\label{lem:takeup}
	Sei das Flötenprofil gemäß Definition~\ref{def:flute} durch
	$y(x) = A\sin(2\pi x/\lambda)$ gegeben. Dann ist der
	\emph{Take-up-Faktor} $f = L/\lambda$ -- das Verhältnis der Bogenlänge $L$
	einer Periode zur Periodenlänge $\lambda$ -- für $A/\lambda \ll 1$ gegeben
	durch die Näherung
	\[
		f \;\approx\; 1 + \pi^2 \left(\frac{A}{\lambda}\right)^2.
	\]
\end{lemma}

\begin{proof}
	Die Bogenlänge einer Periode ist
	\[
		L = \int_0^\lambda \sqrt{1 + y'(x)^2}\,dx, \qquad
		y'(x) = A \cdot \frac{2\pi}{\lambda}\cos\!\left(\frac{2\pi x}{\lambda}\right).
	\]
	Für $|y'(x)| \ll 1$ gilt die Taylor-Approximation
	$\sqrt{1+u^2} \approx 1 + \tfrac{1}{2}u^2$. Damit
	\[
		L \approx \int_0^\lambda \left(1 + \tfrac{1}{2}y'(x)^2\right) dx
		= \lambda + \frac{1}{2}\int_0^\lambda y'(x)^2\,dx.
	\]
	Das verbleibende Integral berechnet sich mittels der Identität
	$\cos^2(\theta) = \tfrac{1}{2}(1+\cos(2\theta))$:
	\[
		\int_0^\lambda y'(x)^2\,dx
		= A^2 \left(\frac{2\pi}{\lambda}\right)^2
		\int_0^\lambda \cos^2\!\left(\frac{2\pi x}{\lambda}\right) dx
		= A^2 \left(\frac{2\pi}{\lambda}\right)^2 \cdot \frac{\lambda}{2}
		= \frac{2\pi^2 A^2}{\lambda}.
	\]
	Einsetzen liefert
	$L \approx \lambda + \dfrac{\pi^2 A^2}{\lambda}$, also
	\[
		f = \frac{L}{\lambda} \approx 1 + \pi^2\left(\frac{A}{\lambda}\right)^2.
	\]
\end{proof}

\begin{bemerkung}
	Der Take-up-Faktor $f$ gibt an, um welchen Faktor mehr Papier für die
	Mittellage benötigt wird, verglichen mit einer flachen Bahn gleicher
	projizierter Länge $\lambda$. Für reale Flötentypen liegt
	$A/\lambda$ typischerweise im Bereich $[0{,}15;\,0{,}3]$ (siehe
	Beispiel~\ref{bsp:fluten}), sodass $f \in [1{,}2;\,1{,}9]$ -- d.\,h. die
	Mittellage benötigt $20\,\%$ bis $90\,\%$ mehr Papier als die Decklagen pro
	Flächeneinheit. Wie in Abbildung~3 ersichtlich, unterschätzt die
	Näherung aus Lemma~\ref{lem:takeup} den exakten Wert systematisch leicht,
	bleibt aber für $A/\lambda \leq 0{,}3$ eine brauchbare Abschätzung mit einem
	relativen Fehler von unter $10\,\%$.
\end{bemerkung}

\begin{beispiel}[Flötentypen]
	\label{bsp:fluten}
	Die Wellpappenindustrie unterscheidet genormte Flötentypen, die sich in
	Höhe $d = 2A$ und Teilung $\lambda$ unterscheiden:

	\begin{center}
	\begin{tabular}{lccc}
		\textbf{Flötentyp} & \textbf{Höhe $d$ (mm)} & \textbf{Teilung $\lambda$ (mm)} & \textbf{$A/\lambda$} \\
		\hline
		A & 4{,}7 & 8{,}5 & 0{,}276 \\
		C & 3{,}6 & 7{,}5 & 0{,}240 \\
		B & 2{,}5 & 6{,}5 & 0{,}192 \\
		E & 1{,}15 & 3{,}5 & 0{,}164 \\
	\end{tabular}
	\end{center}

	Flötentyp A bietet die größte Flötenhöhe und damit (nach
	Satz~\ref{satz:steifigkeit}) das größte Trägheitsmoment, benötigt jedoch
	auch den größten Take-up-Faktor und damit den höchsten relativen
	Papierverbrauch. Flötentyp E ist materialsparend, aber strukturell
	schwächer -- ein erstes Beispiel für den in Abschnitt~4 formalisierten
	Zielkonflikt.
\end{beispiel}

\subsection{Tragfähigkeitsmodell}

\begin{satz}[Trägheitsmoment des Sandwichquerschnitts]
	\label{satz:traegheit}
	Ein Wandelement der Wellpappe mit Breite $b$ bestehe aus zwei Decklagen der
	Dicke $t_\ell$, deren Mittelflächen den Abstand $d$ (Flötenhöhe) haben. Das
	Flächenträgheitsmoment $I_{\mathrm{sand}}$ des Querschnitts bezüglich der
	neutralen Faser (Mittelebene) ist
	\[
		I_{\mathrm{sand}} = \frac{b\,t_\ell^3}{6} + \frac{b\,t_\ell\,d^2}{2}.
	\]
\end{satz}

\begin{proof}
	Jede der beiden Decklagen ist ein Rechteck mit Breite $b$ und Höhe
	$t_\ell$, dessen eigenes Flächenträgheitsmoment bezüglich des eigenen
	Schwerpunkts $\dfrac{b\,t_\ell^3}{12}$ beträgt. Der Schwerpunkt jeder
	Decklage liegt im Abstand $d/2$ von der neutralen Faser. Nach dem
	Satz von Steiner (Parallelachsensatz) ist das Trägheitsmoment einer
	einzelnen Decklage bezüglich der neutralen Faser
	\[
		I_{\text{einzeln}} = \frac{b\,t_\ell^3}{12} + b\,t_\ell \left(\frac{d}{2}\right)^2
		= \frac{b\,t_\ell^3}{12} + \frac{b\,t_\ell\,d^2}{4}.
	\]
	Da der Querschnitt aus zwei symmetrisch angeordneten Decklagen besteht
	(der Beitrag der gewellten Mittellage selbst wird für die Biegesteifigkeit
	als vernachlässigbar angenommen, da sie nur über die Stützwirkung wirkt),
	ergibt sich das Gesamtträgheitsmoment durch Verdopplung:
	\[
		I_{\mathrm{sand}} = 2\,I_{\text{einzeln}}
		= \frac{b\,t_\ell^3}{6} + \frac{b\,t_\ell\,d^2}{2}. \qedhere
	\]
\end{proof}

\begin{satz}[Steifigkeitsgewinn durch Sandwichstruktur]
	\label{satz:steifigkeit}
	Sei $I_{\mathrm{flach}} = \dfrac{b(2t_\ell)^3}{12} = \dfrac{2}{3}b\,t_\ell^3$
	das Flächenträgheitsmoment einer massiven, flachen Bahn mit derselben
	Gesamt-Decklagenmasse (Dicke $2t_\ell$, gleiches $b$). Dann gilt für
	$t_\ell \ll d$:
	\[
		\frac{I_{\mathrm{sand}}}{I_{\mathrm{flach}}}
		\;\approx\; \frac{3}{4}\left(\frac{d}{t_\ell}\right)^2.
	\]
\end{satz}

\begin{proof}
	Für $t_\ell \ll d$ ist der Term $\dfrac{b\,t_\ell^3}{6}$ in
	$I_{\mathrm{sand}}$ gegenüber $\dfrac{b\,t_\ell\,d^2}{2}$ von
	höherer Ordnung klein und kann vernachlässigt werden:
	$I_{\mathrm{sand}} \approx \dfrac{b\,t_\ell\,d^2}{2}$. Damit
	\[
		\frac{I_{\mathrm{sand}}}{I_{\mathrm{flach}}}
		\approx \frac{\frac{1}{2}b\,t_\ell\,d^2}{\frac{2}{3}b\,t_\ell^3}
		= \frac{1}{2} \cdot \frac{3}{2} \cdot \frac{d^2}{t_\ell^2}
		= \frac{3}{4}\left(\frac{d}{t_\ell}\right)^2. \qedhere
	\]
\end{proof}

\begin{korollar}
	Für ein typisches Verhältnis $d/t_\ell = 20$ (z.\,B. $d \approx 4\,\mathrm{mm}$,
	$t_\ell \approx 0{,}2\,\mathrm{mm}$) ergibt sich
	$I_{\mathrm{sand}}/I_{\mathrm{flach}} \approx \tfrac{3}{4}\cdot 400 = 300$.
	Die Wellpappenstruktur erreicht also bei \emph{identischer Papiermasse} ein
	etwa $300$-fach größeres Flächenträgheitsmoment als eine massive, flache
	Bahn -- dies ist der zentrale physikalische Mechanismus, der die
	außergewöhnliche Steifigkeit von Wellpappe bei minimalem Materialeinsatz
	erklärt (siehe Abbildung~5).
\end{korollar}

\begin{satz}[Euler-Knicklast einer Kartonwand]
	\label{satz:euler}
	Eine Kartonwand der Höhe $L$ (Stapelrichtung), modelliert als schlanke,
	beidseitig gelenkig gelagerte Platte mit Biegesteifigkeit $E I_{\mathrm{sand}}$
	($E$: Elastizitätsmodul des Papiers, $I_{\mathrm{sand}}$ gemäß
	Satz~\ref{satz:traegheit}), besitzt die kritische Knicklast
	\[
		P_{\mathrm{cr}} = \frac{\pi^2 E\,I_{\mathrm{sand}}}{L^2}.
	\]
\end{satz}

\begin{proof}
	Die Differentialgleichung der Biegelinie eines axial belasteten,
	gelenkig gelagerten Stabes lautet
	\[
		E I_{\mathrm{sand}}\,y''(x) + P\,y(x) = 0, \qquad y(0) = y(L) = 0.
	\]
	Mit $k^2 = \dfrac{P}{E I_{\mathrm{sand}}}$ lautet die allgemeine Lösung
	$y(x) = C_1 \sin(kx) + C_2\cos(kx)$. Die Randbedingung $y(0)=0$ liefert
	$C_2 = 0$. Die Randbedingung $y(L)=0$ liefert $C_1 \sin(kL) = 0$. Für eine
	nichttriviale Lösung ($C_1 \neq 0$) muss $\sin(kL) = 0$ gelten, also
	$kL = n\pi$ für $n \in \mathbb{N}$. Die kleinste positive Lösung ($n=1$,
	die erste und damit praktisch relevante Knickform) ergibt $k = \pi/L$,
	also
	\[
		P_{\mathrm{cr}} = k^2 E I_{\mathrm{sand}} = \frac{\pi^2 E\,I_{\mathrm{sand}}}{L^2}. \qedhere
	\]
\end{proof}

\begin{bemerkung}[McKee-Formel]
	\label{bem:mckee}
	Die in Satz~\ref{satz:euler} hergeleitete Knicklast beschreibt das
	Versagen einer \emph{einzelnen, idealisierten} Wand. Die tatsächliche
	Stapeldruckfestigkeit eines vollständigen Kartons hängt zusätzlich vom
	lokalen Kantenversagen (Edge Crush) sowie von Wechselwirkungen zwischen
	den vier Seitenwänden ab und entzieht sich einer rein analytischen
	Herleitung. McKee, Gander und Wachuta (1963) entwickelten daher die
	semi-empirische Formel
	\[
		\mathrm{BCT} \;=\; 5{,}87 \cdot \mathrm{ECT}^{0{,}746}
		\cdot h^{0{,}508} \cdot Z^{0{,}492},
	\]
	wobei $\mathrm{ECT}$ der Kantenstauchwiderstand, $h$ das Kaliber
	(Gesamtdicke) der Wellpappe und $Z = 2(\ell+b)$ der Umfang des Kartons ist.
	Die Exponenten von $h$ und $Z$ liegen nahe bei den Werten $1$ bzw.\ $-1$,
	die eine reine Plattenbeul-Theorie (analog Satz~\ref{satz:euler}, bei der
	$P_{\mathrm{cr}} \propto I_{\mathrm{sand}} \propto t_\ell d^2$ und
	$P_{\mathrm{cr}} \propto 1/L^2$ gilt) nahelegen würde; die empirisch
	bestimmten, etwas kleineren Exponenten ($0{,}508$ statt $1$, $0{,}492$
	statt $1$, und insbesondere ein \emph{positiver} statt negativer Exponent
	für $Z$) reflektieren den Übergang vom reinen Knickversagen zu einem
	kombinierten Versagensmodus mit lokalem Kantenstauchen, der bei realen,
	endlich steifen Kartons dominiert. Die McKee-Formel ist bis heute der
	industrielle Standard zur Auslegung von Versandkartons (vgl.
	Abbildung~4).
\end{bemerkung}

\section{Analyse}

\subsection{Gesamtoptimierung als Mehrzieloptimierung}

Die Sätze~\ref{satz:wuerfel} und \ref{satz:offen} minimieren den Materialeinsatz
\emph{ohne} Berücksichtigung der Tragfähigkeit, während Satz~\ref{satz:steifigkeit}
und Bemerkung~\ref{bem:mckee} die Tragfähigkeit beschreiben, ohne den
Materialeinsatz zu begrenzen. Das vollständige Problem ist ein
Mehrzieloptimierungsproblem:

\begin{satz}[Optimalitätsbedingung unter Tragfähigkeitsrestriktion]
	\label{satz:lagrange}
	Sei $V > 0$ fest und $F_{\min} > 0$ eine geforderte Mindeststapellast.
	Betrachte das Optimierungsproblem
	\[
		\min_{\ell,b,h > 0} S(\ell,b,h) \quad \text{u.\,d.\,N.} \quad
		\ell b h = V \;\;\text{und}\;\; \mathrm{BCT}(\ell,b,h) \geq F_{\min}.
	\]
	Ist die BCT-Restriktion am Optimum aktiv (d.\,h.
	$\mathrm{BCT}(\ell^\ast,b^\ast,h^\ast) = F_{\min}$), so existieren
	Lagrange-Multiplikatoren $\mu, \nu \in \mathbb{R}$ mit
	\[
		\nabla S(\ell^\ast,b^\ast,h^\ast)
		= \mu \,\nabla(\ell b h)\big|_{(\ell^\ast,b^\ast,h^\ast)}
		+ \nu \,\nabla \mathrm{BCT}(\ell^\ast,b^\ast,h^\ast),
		\qquad \nu \geq 0.
	\]
\end{satz}

\begin{proof}
	Dies ist eine direkte Anwendung der Karush-Kuhn-Tucker-Bedingungen (KKT)
	für ein Optimierungsproblem mit einer Gleichheits- und einer
	Ungleichheitsrestriktion. Die Zielfunktion $S$ und die Restriktionsfunktion
	$\mathrm{BCT}$ sind auf dem zulässigen Bereich $\ell,b,h>0$ stetig
	differenzierbar (als Komposition von Potenz- und Polynomfunktionen), und
	der zulässige Bereich ist nichtleer und offen, sodass eine
	Regularitätsbedingung (LICQ) generisch erfüllt ist. Die KKT-Theorie
	garantiert dann, dass an einem lokalen Minimum $(\ell^\ast,b^\ast,h^\ast)$
	der Gradient der Zielfunktion eine nichtnegative Linearkombination der
	Gradienten der aktiven Restriktionen ist (mit beliebigem Vorzeichen für die
	Gleichheitsrestriktion und nichtnegativem Vorzeichen $\nu \geq 0$ für die
	aktive Ungleichheitsrestriktion, da diese in Richtung wachsendem
	$\mathrm{BCT}$ zeigt und $S$ am Optimum nicht durch eine Verbesserung von
	$\mathrm{BCT}$ weiter reduziert werden kann).
\end{proof}

\begin{bemerkung}
	Satz~\ref{satz:lagrange} formalisiert den intuitiven Befund, dass das
	optimale Seitenverhältnis eines Kartons \emph{von der geforderten
	Tragfähigkeit abhängt}: Ist $F_{\min}$ klein (leichte Güter), so ist die
	BCT-Restriktion inaktiv ($\nu = 0$) und Satz~\ref{satz:wuerfel} bzw.
	Satz~\ref{satz:offen} liefert die Lösung (würfelnahe Form). Ist $F_{\min}$
	groß (schwere, stapelbare Güter), so wird die BCT-Restriktion aktiv und das
	Optimum verschiebt sich hin zu Formen mit größerem Umfang $Z$ (gemäß
	Bemerkung~\ref{bem:mckee} steigt $\mathrm{BCT}$ mit $Z^{0{,}492}$), selbst
	wenn dies die Oberfläche $S$ erhöht. Die beobachtete Praxis -- Versandkartons
	für schwere Güter sind tendenziell flacher und breiter als rein
	volumenoptimierte Würfel -- ist damit eine direkte Konsequenz der hier
	formal hergeleiteten Optimalitätsbedingung.
\end{bemerkung}

\subsection{Sicherheitsfaktoren und maximale Stapelhöhe}

\begin{beispiel}[Maximale Stapelhöhe]
	Sei $\mathrm{BCT} = 4500\,\mathrm{N}$ (typischer Wert für einen mittelgroßen
	Versandkarton mit $\mathrm{ECT} = 6\,\mathrm{kN/m}$, $h = 4\,\mathrm{mm}$,
	$Z = 140\,\mathrm{cm}$, vgl. Abbildung~4) und das Gewicht eines gefüllten
	Kartons $m = 12\,\mathrm{kg}$, also $G = mg \approx 117{,}7\,\mathrm{N}$
	(mit $g = 9{,}81\,\mathrm{m/s^2}$). In der Praxis wird ein
	\emph{Sicherheitsfaktor} $\mathrm{SF} \in [4,6]$ angesetzt, um Effekte wie
	Feuchtigkeit, Lagerzeit (Kriechen des Papiers) und dynamische Stöße beim
	Transport zu berücksichtigen:
	\[
		\mathrm{BCT}_{\mathrm{zul}} = \frac{\mathrm{BCT}}{\mathrm{SF}}.
	\]
	Die maximale Anzahl $n_{\max}$ stapelbarer Kartons (der unterste Karton
	trägt das Gewicht aller $n-1$ darüberliegenden) ergibt sich aus
	$(n_{\max}-1) \cdot G \leq \mathrm{BCT}_{\mathrm{zul}}$, also
	\[
		n_{\max} = \left\lfloor \frac{\mathrm{BCT}}{\mathrm{SF}\cdot G} \right\rfloor + 1.
	\]
	Für $\mathrm{SF} = 5$ ergibt sich
	$n_{\max} = \left\lfloor \dfrac{4500}{5 \cdot 117{,}7} \right\rfloor + 1
	= \lfloor 7{,}64 \rfloor + 1 = 8$. Ein einzelner Stapel sollte also aus
	höchstens $8$ solchen Kartons bestehen -- eine in der Logistikpraxis
	unmittelbar verwertbare Konsequenz der zuvor rein formal hergeleiteten
	Gleichungen.
\end{beispiel}

\subsection{Vergleich mit Alternativmaterialien}

\begin{bemerkung}
	Die in Satz~\ref{satz:steifigkeit} quantifizierte Steifigkeits-Masse-Relation
	erklärt auch, warum Wellpappe gegenüber Alternativen wie Massivholz,
	Vollkunststoff oder geschäumten Polymeren in vielen Anwendungen überlegen
	ist: Massivholz erreicht zwar ein hohes $E$, aber bei deutlich höherer
	Dichte und ohne den quadratischen Hebel des Abstandsterms $d^2$ in
	$I_{\mathrm{sand}}$. Geschäumte Polymere (z.\,B. EPS) erreichen einen
	ähnlichen Sandwich-Effekt durch eingeschlossene Luft, sind jedoch
	energieintensiver in der Herstellung und schwerer recycelbar. Die
	Kombination aus zellulosebasiertem Rohstoff (nachwachsend, recycelbar),
	geometrischer Versteifung durch Sinusprofil (Satz~\ref{satz:steifigkeit})
	und global optimierter Quaderform (Sätze~\ref{satz:wuerfel}
	und~\ref{satz:offen}) macht den Karton somit nicht nur mathematisch,
	sondern auch ökologisch und ökonomisch zu einer herausragenden Lösung.
\end{bemerkung}

\section{Numerische Veranschaulichung}

Die folgenden Abbildungen veranschaulichen die in den vorangegangenen Abschnitten
hergeleiteten Ergebnisse anhand numerischer Beispiele (erzeugt mit
\texttt{matplotlib}).

\begin{figure}[h]
	\centering
	\includegraphics[width=0.85\textwidth]{plots/abb1_materialbedarf.pdf}
	\caption{Materialbedarf (Oberfläche) $S(\ell)$ für geschlossene und offene
	Kartons mit quadratischer Grundfläche und festem Volumen $V=1$, in
	Abhängigkeit der Kantenlänge $\ell$. Die markierten Minima entsprechen den
	in Satz~\ref{satz:wuerfel} (geschlossen, $\ell^\ast = V^{1/3}$) und
	Satz~\ref{satz:offen} (offen, $\ell^\ast = (2V)^{1/3}$) hergeleiteten
	optimalen Kantenlängen. Sichtbar ist, dass der offene Karton bei seinem
	jeweiligen Optimum eine größere Grundfläche, aber eine insgesamt kleinere
	Oberfläche besitzt als der geschlossene Karton bei seinem Optimum.}
\end{figure}

\begin{figure}[h]
	\centering
	\includegraphics[width=0.85\textwidth]{plots/abb2_flutenprofile.pdf}
	\caption{Sinusmodell (Definition~\ref{def:flute}) der Flötenprofile für die
	vier in Beispiel~\ref{bsp:fluten} aufgeführten Standard-Wellenarten A, B, C
	und E. Deutlich erkennbar ist, dass Flötentyp A die größte Amplitude und die
	größte Periodenlänge besitzt, während Flötentyp E die feinste und
	flachste Welle aufweist. Die Fläche unter den Kurven (im Sinne der
	Bogenlänge) bestimmt direkt den in Lemma~\ref{lem:takeup} behandelten
	Take-up-Faktor.}
\end{figure}

\begin{figure}[h]
	\centering
	\includegraphics[width=0.85\textwidth]{plots/abb3_takeup_faktor.pdf}
	\caption{Take-up-Faktor $f = L/\lambda$ als Funktion des Verhältnisses
	$A/\lambda$, sowohl exakt (numerische Integration der Bogenlänge) als auch
	gemäß der Näherungsformel aus Lemma~\ref{lem:takeup}. Die vier
	Standard-Flötentypen sind als Punkte markiert. Für die in der Praxis
	relevanten Werte $A/\lambda \leq 0{,}3$ approximiert die Formel
	$f \approx 1 + \pi^2(A/\lambda)^2$ den exakten Wert mit einer relativen
	Abweichung von unter $10\,\%$; für größere Verhältnisse divergiert die
	Näherung zunehmend vom exakten Integral, da die zugrundeliegende
	Taylor-Approximation $\sqrt{1+u^2}\approx 1+\tfrac{1}{2}u^2$ nur für kleine
	$u$ gültig ist.}
\end{figure}

\begin{figure}[h]
	\centering
	\includegraphics[width=0.85\textwidth]{plots/abb4_mckee_formel.pdf}
	\caption{Stapeldruckfestigkeit $\mathrm{BCT}$ nach der McKee-Formel
	(Bemerkung~\ref{bem:mckee}) als Funktion des Kartonumfangs
	$Z = 2(\ell+b)$ für ein festes Kaliber $h = 4\,\mathrm{mm}$ und vier
	verschiedene Kantenstauchwiderstände $\mathrm{ECT} \in \{3,5,7,9\}\,
	\mathrm{kN/m}$. Der sublineare Anstieg ($\mathrm{BCT} \propto
	Z^{0{,}492}$) zeigt, dass eine Vergrößerung des Kartons die Tragfähigkeit
	zwar erhöht, jedoch mit abnehmendem Grenzertrag -- eine Verdopplung des
	Umfangs erhöht die Stapeldruckfestigkeit nur um den Faktor
	$2^{0{,}492}\approx 1{,}406$.}
\end{figure}

\begin{figure}[h]
	\centering
	\includegraphics[width=0.85\textwidth]{plots/abb5_traegheitsmoment.pdf}
	\caption{Verhältnis der Flächenträgheitsmomente
	$I_{\mathrm{sand}}/I_{\mathrm{flach}}$ als Funktion des Verhältnisses
	$d/t_\ell$ (Flötenhöhe zu Decklagendicke), gemäß
	Satz~\ref{satz:steifigkeit}. Der quadratische Zusammenhang
	$I_{\mathrm{sand}}/I_{\mathrm{flach}} \approx \tfrac{3}{4}(d/t_\ell)^2$
	zeigt, dass bereits moderate Vergrößerungen des Flötenhöhe-zu-Decklagendicke-
	Verhältnisses zu drastischen Steifigkeitsgewinnen führen. Für das
	praxisrelevante Verhältnis $d/t_\ell = 20$ (markiert) ergibt sich ein
	Faktor von ca.\ $300$ -- der zentrale quantitative Beleg für die
	„Genialität'' der Wellpappenstruktur.}
\end{figure}

\section{Zusammenfassung}

Diese Arbeit hat den Karton -- konkret die Wellpappenverpackung -- als Lösung zweier
gekoppelter, formal lösbarer Optimierungsprobleme analysiert:

\begin{itemize}
	\item Auf der \textbf{Makroebene} minimiert die Quaderform bei gegebenem
	      Volumen den Materialeinsatz (Sätze~\ref{satz:wuerfel}
	      und~\ref{satz:offen}, mit dem charakteristischen
	      $V^{2/3}$-Skalierungsgesetz).
	\item Auf der \textbf{Mikroebene} maximiert das sinusförmige Flötenprofil
	      bei gegebener Papiermasse das Flächenträgheitsmoment und damit die
	      Knicksteifigkeit (Sätze~\ref{satz:traegheit}--\ref{satz:euler}, mit
	      dem quadratischen Verstärkungsfaktor $\tfrac{3}{4}(d/t_\ell)^2$).
	\item Die \textbf{Kopplung} beider Ebenen unter einer
	      Tragfähigkeitsrestriktion führt zu einer Lagrange-Charakterisierung
	      (Satz~\ref{satz:lagrange}), die das in der Praxis beobachtete
	      Verhalten -- größere, flachere Kartons für schwerere Güter --
	      formal erklärt.
\end{itemize}

\subsection{Anwendungen}

Die hergeleiteten Modelle besitzen unmittelbare praktische Anwendungen:

\begin{itemize}
	\item \textbf{Verpackungsdesign}: Automatisierte Bestimmung optimaler
	      Kartonabmessungen und Flötentypen für ein gegebenes
	      Produktportfolio, unter Einhaltung von Tragfähigkeitsvorgaben.
	\item \textbf{Logistikplanung}: Berechnung maximaler Stapelhöhen
	      (Abschnitt~4.2) zur sicheren Lagerplanung in Hochregallagern.
	\item \textbf{Materialeinkauf}: Quantifizierung des Take-up-Faktors
	      (Lemma~\ref{lem:takeup}) zur exakten Kalkulation des
	      Papierbedarfs verschiedener Flötentypen.
	\item \textbf{Qualitätssicherung}: Nutzung der McKee-Formel
	      (Bemerkung~\ref{bem:mckee}) zur Spezifikation und Prüfung von
	      Mindestanforderungen an Lieferkartons.
	\item \textbf{Bildung und Popularisierung}: Demonstration, wie
	      elementare Werkzeuge der Analysis, Differentialgeometrie und
	      Balkentheorie auf ein allgegenwärtiges Alltagsobjekt angewendet
	      werden können -- ein anschauliches Beispiel für angewandte
	      Mathematik im Ingenieurwesen.
\end{itemize}

\subsection{Ausblick}

Die in dieser Arbeit entwickelten Modelle bilden lediglich den Ausgangspunkt für
eine Vielzahl vielversprechender Erweiterungen, die im Folgenden ausführlich
diskutiert werden.

\paragraph{Topologieoptimierung und generatives Design.}
Die hier betrachteten Modelle gehen von einer festen Quaderform und einem
einfachen Sinusprofil aus. Moderne Verfahren der Topologieoptimierung, kombiniert
mit der Finite-Elemente-Methode (FEM), könnten genutzt werden, um nicht-sinusförmige,
lokal variierende Flötenprofile zu entwerfen, die das Verhältnis
$I_{\mathrm{sand}}/I_{\mathrm{flach}}$ aus Satz~\ref{satz:steifigkeit} weiter
verbessern -- etwa durch trapezförmige, asymmetrische oder fraktale Wellenmuster,
deren Trägheitsmomente analytisch deutlich aufwendiger, aber numerisch gut
zugänglich wären. Generative Designalgorithmen könnten zudem die Quaderform selbst
aufgeben und organische, an die Produktgeometrie angepasste Kartonformen
erzeugen, deren Materialbedarf das $V^{2/3}$-Skalierungsgesetz (Bemerkung nach
Satz~\ref{satz:offen}) gegebenenfalls unterschreitet, sofern die
Tragfähigkeitsrestriktion entsprechend angepasst wird.

\paragraph{Hygroskopisches Verhalten und Klimaeinflüsse.}
Das in Satz~\ref{satz:euler} verwendete Elastizitätsmodul $E$ von Papier ist
stark feuchtigkeitsabhängig: Bei hoher relativer Luftfeuchtigkeit kann $E$ um
$30$--$50\,\%$ sinken, was über $P_{\mathrm{cr}} \propto E$ direkt die Knicklast
reduziert. Eine Erweiterung des Modells um eine explizite Abhängigkeit
$E = E(\varphi, t)$ (relative Feuchtigkeit $\varphi$, Lagerzeit $t$, unter
Berücksichtigung von Kriechprozessen im Papiergefüge) würde die
Sicherheitsfaktoren aus Abschnitt~4.2 nicht mehr als pauschale Konstanten,
sondern als Funktionen der Lager- und Transportbedingungen formalisieren --
ein wichtiger Schritt hin zu klimaadaptiven Verpackungsspezifikationen,
insbesondere für den globalen Schiffstransport durch Klimazonen mit stark
unterschiedlicher Luftfeuchtigkeit.

\paragraph{Aktive und intelligente Verpackungen (Smart Packaging).}
Zukünftige Kartons könnten mit eingebetteten Sensoren (Feuchtigkeit,
Temperatur, Beschleunigung, Neigung) sowie RFID- oder Bluetooth-Trackern
ausgestattet werden. Die hier entwickelten formalen Modelle ließen sich dann zu
\emph{dynamischen} Modellen erweitern: Die Stapeldruckfestigkeit
$\mathrm{BCT}(t)$ als Funktion der Zeit könnte mittels der gemessenen
Feuchtigkeits- und Lasthistorie laufend neu geschätzt werden, sodass
Logistikunternehmen in Echtzeit vor drohendem Versagen einzelner Stapel gewarnt
werden könnten. Dies erfordert eine Kopplung der mechanischen Modelle dieser
Arbeit mit Methoden des maschinellen Lernens zur Zeitreihenprognose.

\paragraph{Mehrzieloptimierung mittels evolutionärer Algorithmen.}
Satz~\ref{satz:lagrange} liefert notwendige Optimalitätsbedingungen für den Fall
einer aktiven Restriktion, jedoch keine geschlossene Lösung für das allgemeine
Mehrzieloptimierungsproblem (Minimierung von Material \emph{und} Maximierung von
Tragfähigkeit \emph{und} Minimierung des Take-up-Faktors). Eine naheliegende
Erweiterung besteht in der numerischen Bestimmung der vollständigen
Pareto-Front mittels genetischer Algorithmen oder Partikelschwarm-Optimierung,
unter Einbeziehung diskreter Entscheidungsvariablen (Flötentyp, Decklagentyp,
Anzahl der Wellpappenlagen: einwellig, zweiwellig, dreiwellig).

\paragraph{Verbindung zum Subgraph Algorithmus: Nesting- und Cutting-Stock-Probleme.}
Ein methodisch interessanter Anknüpfungspunkt ergibt sich aus der Perspektive
der Schnittmusteroptimierung (\emph{Cutting-Stock-Problem}): Bevor ein Karton
gefaltet wird, liegt er als zweidimensionaler Stanzzuschnitt (\emph{die-cut})
vor, dessen Faltlinien und Klebelaschen eine strukturierte, graphartige Topologie
besitzen (Knoten = Falzkanten, Kanten = angrenzende Flächen). Die effiziente
Anordnung (Nesting) mehrerer solcher Zuschnitte auf einer Wellpappenbahn mit
minimalem Verschnitt lässt sich als Subgraph-Erkennungsproblem auffassen:
Gesucht werden wiederkehrende strukturelle Muster zwischen verschiedenen
Zuschnitt-Layouts, analog zur in früheren Arbeiten beschriebenen Anwendung des
Subgraph Algorithmus auf die Analyse von Abstract Syntax Trees. Eine
Adjazenzmatrix-Repräsentation der Faltlinien-Topologie, kombiniert mit den
Signatur-Arrays und der zyklischen Rotation des Subgraph Algorithmus, könnte
verwendet werden, um modulare Stanzformen zu identifizieren, die sich mit
minimalem Verschnitt zu größeren Bahnen kombinieren lassen -- eine direkte
Brücke zwischen der hier vorgestellten Kontinuumsmechanik des Kartons und der
diskreten Graphentheorie.

\paragraph{Automatisierte Faltrobotik und Verpackungslinien.}
Die geometrischen Optimalitätsbedingungen aus Sätzen~\ref{satz:wuerfel}
und~\ref{satz:offen} setzen voraus, dass beliebige Kantenlängen $\ell, b, h$
realisierbar sind. In automatisierten Verpackungslinien ist dies jedoch durch
die Mechanik der Faltroboter (Greifarme, Faltschwerter, Klebedüsen)
eingeschränkt, die nur diskrete Formfaktoren effizient verarbeiten können. Eine
Erweiterung des Modells um diskrete Formfaktor-Restriktionen (z.\,B.
$\ell, b, h$ aus einem vorgegebenen Raster) und um die
Zykluszeiten der Faltroboter als zusätzliches Optimierungsziel würde das Modell
für die industrielle Praxis direkt einsetzbar machen.

\paragraph{Kreislaufwirtschaft, Recyclingfasern und Lebenszyklusanalyse.}
Das Elastizitätsmodul $E$ und der Kantenstauchwiderstand $\mathrm{ECT}$ hängen
vom Anteil an Recyclingfasern im Papier ab: Mit jedem Recyclingzyklus verkürzen
sich die Zellulosefasern, was $E$ und $\mathrm{ECT}$ geringfügig reduziert. Eine
Erweiterung der McKee-Formel (Bemerkung~\ref{bem:mckee}) um einen
Alterungsparameter $\mathrm{ECT}(n_{\text{Recycling}})$ würde es erlauben, die
in dieser Arbeit hergeleiteten Optimalitätsbedingungen mit einer vollständigen
Lebenszyklusanalyse (LCA) zu koppeln: Optimierung nicht nur des
Materialeinsatzes pro Karton, sondern des gesamten ökologischen Fußabdrucks über
mehrere Recyclingzyklen hinweg, inklusive Energieaufwand für Wiederaufbereitung
und Transport.

\paragraph{Erweiterung auf gewichtete, dynamische und stoßartige Belastungen.}
Die Sätze~\ref{satz:euler} und die McKee-Formel (Bemerkung~\ref{bem:mckee})
betrachten ausschließlich statische, vertikale Druckbelastung. In der Realität
sind Kartons jedoch dynamischen Stoßbelastungen (Fallprüfung, Vibrationen
während des Transports) sowie horizontalen Scherkräften ausgesetzt. Eine
Erweiterung um ein dynamisches Modell -- etwa basierend auf einem
Feder-Dämpfer-System mit nichtlinearer Steifigkeit des Wellpappenkerns --
würde es erlauben, die Fallhöhen-Toleranz eines Kartons formal aus den hier
entwickelten geometrischen und materialspezifischen Größen ($d$, $t_\ell$,
$\lambda$, $E$) herzuleiten, anstatt sie empirisch durch Falltests zu
bestimmen.

\paragraph{Digitale Zwillinge für Verpackungssysteme.}
Die Kombination aller oben genannten Erweiterungen -- feuchtigkeitsabhängige
Materialparameter, Sensordaten, dynamische Belastungsmodelle und
Lebenszyklusdaten -- legt die Entwicklung eines vollständigen \emph{digitalen
Zwillings} eines Kartons nahe: ein Software-Modell, das für jeden individuellen
Karton (identifizierbar über eingebettete RFID-Tags) kontinuierlich seinen
aktuellen Tragfähigkeitszustand $\mathrm{BCT}(t)$ auf Basis der formalen Modelle
dieser Arbeit, kombiniert mit Echtzeit-Sensordaten, berechnet und damit
prädiktive Aussagen über die Resttragfähigkeit eines Stapels über die gesamte
Lieferkette hinweg ermöglicht.

\paragraph{Verallgemeinerung auf nicht-quaderförmige und additiv gefertigte Strukturen.}
Schließlich beschränkt sich diese Arbeit auf quaderförmige Kartons mit
sinusförmigem Flötenprofil. Mit dem Aufkommen des 3D-Drucks von
Wellpappenstrukturen und faserbasierten Verbundmaterialien wird es zunehmend
möglich, beliebige, lokal variierende Wandstrukturen zu fertigen. Eine
Verallgemeinerung der hier entwickelten Theorie auf variable Querschnittsprofile
$y(x)$ (statt der festen Sinusform aus Definition~\ref{def:flute}) und auf
nicht-quaderförmige Hüllgeometrien (z.\,B. polyedrische oder zylindrische
Verpackungen) würde den Anwendungsbereich der hier formal etablierten
Optimalitätsprinzipien -- Minimierung der Oberfläche bei gegebenem Volumen
(Makroebene) und Maximierung des Trägheitsmoments bei gegebener Materialmasse
(Mikroebene) -- erheblich erweitern und auf eine neue Klasse von
„algorithmisch entworfenen'' Verpackungsmaterialien übertragen.

\begin{thebibliography}{99}

\bibitem{mckee1963}
McKee, R.~C., Gander, J.~W., Wachuta, J.~R.:
\emph{Compression Strength Formula for Corrugated Boxes}.
Paper Trade Journal, Vol.~149, No.~8, 1963.

\bibitem{maltenfort1956}
Maltenfort, G.~G.:
\emph{Compression Strength of Corrugated Containers}.
Fibre Containers, Vol.~41, No.~7, 1956.

\bibitem{frank2014}
Frank, B.:
\emph{Corrugated Box Compression -- A Literature Survey}.
Packaging Technology and Science, Vol.~27, No.~2, 2014.

\bibitem{twede2005}
Twede, D., Selke, S.:
\emph{Cartons, Crates and Corrugated Board: Handbook of Paper and Wood
Packaging Technology}.
DEStech Publications, Lancaster, PA, 2005.

\bibitem{markstroem1999}
Markström, H.:
\emph{Testing Methods and Standards for Corrugated Board}.
Lorentzen \& Wettre, Stockholm, 1999.

\bibitem{tappi2017}
TAPPI:
\emph{T~804: Compression Test of Fiberboard Shipping Containers}.
TAPPI Test Methods, Atlanta, 2017.

\bibitem{timoshenko}
Timoshenko, S.~P., Gere, J.~M.:
\emph{Theory of Elastic Stability}.
2.~Auflage, McGraw-Hill, New York, 1961.

\bibitem{abboud2015}
Abboud, A., Backurs, A., Vassilevska Williams, V.:
\emph{Tight Hardness Results for LCS and Other Sequence Similarity
Measures}.
In: Proceedings of FOCS 2015.

\bibitem{epp_subgraph}
Epp, S.:
\emph{Der Subgraph Algorithmus}.
GitLab: \url{https://github.com/hjstephan86/subgraph}, 2026.

\end{thebibliography}

\end{document}

